\documentclass[11pt]{article}

% -------------------------------------------------
% PACKAGES
% -------------------------------------------------
\usepackage[letterpaper,margin=1in]{geometry}
\usepackage{amsmath, amssymb, amsthm}
\usepackage{graphicx, xcolor, multirow, enumerate, fancyhdr, caption, listings, wrapfig}
\usepackage{tikz, pgfplots}
\usetikzlibrary{calc, arrows.meta, positioning, decorations.markings, shadings}
\usepackage{booktabs}
\usepackage{authblk} % for multiple authors
\usepackage{lipsum}
\usepackage[none]{hyphenat}

% -------------------------------------------------
% CODE SETTINGS
% -------------------------------------------------
\lstset{
  language=Matlab,
  basicstyle=\ttfamily\small,
  keywordstyle=\color{blue},
  stringstyle=\color{violet},
  commentstyle=\color{teal},
  frame=single,
  framesep=10pt,
  xleftmargin=10pt,
  xrightmargin=10pt,
  breaklines=true,
  numbers=left,
  numberstyle=\tiny\color{gray},
  showstringspaces=false,
  tabsize=4,
  captionpos=b
}

% -------------------------------------------------
% TITLE BLOCK
% -------------------------------------------------
\title{\vspace{-1.2cm} % move title up
\bfseries\Large Multimodal Alpha Forecasting and Robust Portfolio Optimization}

\author[]{Catherine Ma (\texttt{cama12@mit.edu})}
\author[]{Hanna Zhang (\texttt{zeyuzh@mit.edu})}
\affil[]{}

\date{} % hide date completely

% -------------------------------------------------
% DOCUMENT
% -------------------------------------------------
\begin{document}
\maketitle
\vspace{-1.2cm} % move content closer to title

\section{Problem Summary}
Our project aims to develop a multimodal machine learning and optimization framework to forecast short-term stock returns and then construct risk-aware portfolios. Using historical data for S\&P~500 companies, we integrate technical indicators, fundamental ratios, and FinBERT-based sentiment from financial news and corporate filings. The model predicts next-day excess returns, which are then translated into optimal portfolio weights through CVaR optimization, linking predictive modeling with prescriptive decision-making.

\section{Datasets}

We construct a dataset combining market, fundamental, textual, and macro information for S\&P~500 equities. The aim is to capture complementary signals from both quantitative trading activity and qualitative corporate disclosures. Table~\ref{tab:data_summary} in the Appendix lists the datasets and features that we propose to use.



\section{Methods}

\subsection*{Predictive Stage - Next-Day Return Forecasting}
\textbf{(1)}  Construct a panel of daily market indicators (returns, momentum, volatility), quarterly fundamentals (P/E, ROE, leverage, growth), and macro variables (VIX, yield spread).  \textbf{(2)} Apply FinBERT to 10-K/10-Q MD\&A sections and financial news to extract sentiment scores and text embeddings, merged with tabular data. \textbf{(3)} Train XGBoost, XGBoost + FinBERT embeddings, Optimal Regression Tree to predict the stocks' next day return.

\subsection*{Prescriptive Stage - Portfolio Optimization}
Predicted returns feed into a mean-CVaR optimization that maximizes expected return and controls tail risk:
\[
\max_{w}\; \mathbb{E}[r^{\top}w] - \lambda \text{CVaR}_{\alpha}(w),
\]
Weights represent the portfolio decisions allocating capital across assets based on predicted returns.

% \subsection{Evaluation and Backtesting}
% Rolling backtests assess both predictive and portfolio performance using IC, hit rate, Sharpe ratio, Sortino ratio, CVaR, and drawdown.  
% Comparisons between tabular and multimodal models quantify the incremental value of FinBERT-based features.


\section{Challenges}
\textbf{(1) Look-ahead bias:} We will mitigate this by aligning fundamentals and filings strictly by their release dates and applying rolling-window validation for training and testing.  
\textbf{(2) High-dimensional text features:} FinBERT embeddings are computationally intensive; we will reduce dimensionality using PCA and apply regularization to prevent overfitting.  
\textbf{(3) Project scope:} Given time constraint, we may prioritize the predictive modeling stage if portfolio optimization cannot be fully implemented.

\newpage
\appendix

\section{Dataset and Feature Summary}

\begin{table}[htbp]
\centering
\caption{Summary of datasets and key features.}
\label{tab:data_summary}
\renewcommand{\arraystretch}{1.15}
\setlength{\tabcolsep}{4pt}
\begin{tabular}{p{2.7cm} p{3.2cm} p{2.2cm} p{6.5cm}}
\toprule
\textbf{Dataset} & \textbf{Source} & \textbf{Frequency} & \textbf{Key Features / Variables} \\
\midrule

\textbf{Market} &
Yahoo Finance, WRDS (CRSP) &
Daily &
Open, High, Low, Close, Volume (OHLCV); daily return; 5D/20D momentum; 10D/20D realized volatility; MA5/MA20 ratio; RSI; turnover (Volume / Shares Outstanding). \\[0.5em]

\textbf{Fundamentals} &
WRDS (Compustat), AlphaVantage API &
Quarterly $\rightarrow$ Daily &
P/E, P/B, ROE, ROA, gross margin, Debt/Assets, EPS and revenue YoY growth, log(Market Cap). Forward-filled daily after each filing. \\[0.5em]

\textbf{Textual (News)} &
NewsAPI, GDELT, FinBERT &
Daily &
FinBERT sentiment mean and volatility (7D, 30D); sentiment momentum ($\Delta7D$); positive/neutral/negative sentiment ratios; news volume (headline count). \\[0.5em]

\textbf{Textual (Filings)} &
SEC EDGAR, FinBERT &
Quarterly $\rightarrow$ Daily &
FinBERT sentiment of MD\&A and Risk Factors; cosine similarity across filings (tone shift); sentiment change vs previous quarter; PCA-reduced FinBERT embeddings. \\[0.5em]

\textbf{Macro} &
FRED, Nasdaq &
Daily &
VIX, Treasury yield spread (10Y–2Y), PMI, Federal Funds Rate, and S\&P 500 index returns. \\

\bottomrule
\end{tabular}
\end{table}



\end{document}
